\chapter{Unit 2: Exploring Two-Variable Data}

\section{Essential Concepts \& Cheat Sheet}

Two-variable quantitative data is explored using scatterplots and evaluated using direction, form, strength, and unusual points (outliers/influential points).

\begin{formulabox}[Core Formulas: Linear Regression]
\textbf{1. Least-Squares Regression Line (LSRL):}
\[
\hat{y} = a + bx \quad \text{or} \quad \hat{y} = b_0 + b_1 x
\]
where $\hat{y}$ is the predicted response variable and $x$ is the explanatory variable.

\textbf{2. Slope of the Regression Line ($b$):}
\[
b = r \left(\frac{s_y}{s_x}\right)
\]

\textbf{3. $y$-Intercept ($a$):}
\[
a = \bar{y} - b\bar{x}
\]
\textbf{Crucial Property:} The point $(\bar{x}, \bar{y})$ always lies on the regression line.

\textbf{4. Residual ($e$):}
\[
\text{Residual} = \text{Actual } y - \text{Predicted } \hat{y} = y - \hat{y}
\]

\textbf{5. Standard Error of the Residuals ($s$):}
\[
s = \sqrt{\frac{\sum (y_i - \hat{y}_i)^2}{n - 2}} = \sqrt{\frac{\sum e_i^2}{n - 2}}
\]
\end{formulabox}

\section{Correlation ($r$) vs Coefficient of Determination ($r^2$)}

\begin{itemize}
    \item \textbf{Correlation ($r$):} Measures the strength and direction of a \emph{linear} relationship between two quantitative variables. $-1 \le r \le 1$. It has \textbf{no units} and is unaffected by linear scaling or swapping $x$ and $y$.
    \item \textbf{Coefficient of Determination ($r^2$):} The proportion of the variation in the response variable ($y$) that is explained by the linear relationship with the explanatory variable ($x$).
\end{itemize}

\section{Residual Plot Interpretation Rules}

\begin{itemize}
    \item \textbf{Linear Model Appropriate:} Residual plot shows a \textbf{random scatter} of points with no clear pattern around the horizontal residual line $= 0$.
    \item \textbf{Linear Model NOT Appropriate:} Residual plot displays a \textbf{curved pattern} (e.g., U-shape or inverted U-shape) or fan shape (heteroscedasticity).
\end{itemize}

\begin{templatebox}[Score-5 Interpretation Sentence Templates]
\textbf{1. Slope ($b$):} \emph{"For each additional 1 [unit of $x$], the predicted [response variable $y$] increases/decreases by [value of $b$] [units of $y$]."}

\textbf{2. $y$-Intercept ($a$):} \emph{"When the [explanatory variable $x$] is 0 [units of $x$], the predicted [response variable $y$] is [value of $a$] [units of $y$]."}

\textbf{3. $r^2$:} \emph{"Approximately [value of $r^2 \times 100\%$] of the variation in [response variable in context] is accounted for by the linear relationship with [explanatory variable in context]."}

\textbf{4. Standard Error of Residuals ($s$):} \emph{"The typical distance between the actual [response variable in context] and the value predicted by the regression line is approximately [value of $s$] [units of $y$]."}
\end{templatebox}

\begin{trapbox}[AP Exam Trap Alert: Forgetting the "Hat" or "Predicted"]
Always write $\hat{y}$ or explicitly write $\widehat{\text{Score}} = a + b(\text{Study Time})$. Never write $y = a + bx$ without indicating prediction, as graders will deduct points!
\end{trapbox}
